\chapter{Mathematical Notation and Glossary}
\label{apx:notation}

To ensure formal rigor and clarity of exposition, the subsequent
chapters utilize a mathematical notation derived from type theory and
functional semantics. This chapter serves as a quick reference for
the symbology and formal concepts introduced in the formalization of the system.

\section{Notational Conventions}
To ensure formal rigor, consistency, and clarity of exposition across chapters, the formal model adheres to a strict typographic taxonomy comprising five homogeneous symbol families, summarized in Table \ref{tab:notational_styles}.

\begin{table}[htpb]
  \centering
  \small
  \begin{tabular}{|l|l|l|p{5.5cm}|}
    \hline
    \textbf{Conceptual Family} & \textbf{Typographic Style} & \textbf{Primary Symbols} & \textbf{Semantics and Scope} \\ \hline
    \textbf{Analysis Pipeline} & Lowercase Greek & $\phi, \varepsilon, \rho, \gamma$ & Operational transformations and pipeline composition ($\phi = \gamma \circ \rho \circ \varepsilon$). \\ \hline
    \textbf{IR Architecture} & Uppercase Calligraphic & $\mathcal{M}, \mathcal{S}, \mathcal{F}, \mathcal{A}, \mathcal{X}, \mathcal{B}, \mathcal{C}, \mathcal{Q}$ & Structural data model components, containers, and query algebra domain. \\ \hline
    \textbf{Resolution Contexts} & Uppercase Greek \& Calligraphic & $\Gamma, \Delta, \Sigma$ \newline $\mathcal{E}, \mathcal{P}, \mathcal{K}$ & Scoping environments ($\Gamma = \langle \mathcal{E}, \mathcal{P}, \mathcal{K} \rangle$), symbol stack ($\Delta$), and lexical scopes ($\Sigma$). \\ \hline
    \textbf{Semantic Types} & Lowercase Greek & $\iota, \pi, \tau, \sigma, \delta$ & Atomic type references ($\tau$), identifiers ($\iota$), qualified paths ($\pi$), signatures ($\sigma$), and imports ($\delta$). \\ \hline
    \textbf{Structural Members} & Lowercase Latin & $f, p, v, q, k$ & Concrete member instances: fields ($f$), parameters ($p$), variables ($v$), queries ($q$), and type kinds ($k$). \\ \hline
  \end{tabular}
  \caption{Typographic and notational taxonomy adopted across the formal model.}
  \label{tab:notational_styles}
\end{table}

The model relies on the following standard algebraic and set-theoretic conventions:
\begin{itemize}
  \item \textbf{Tuples (Cartesian Product)} $\langle x_1, x_2, \dots,
    x_n \rangle$: Represent composite structures (corresponding to
    \texttt{structs} in the source code). Each element of the tuple
    defines an essential property of the entity.
  \item \textbf{Disjoint Union (Sum Type)} $A \mid B \mid C$:
    Indicates that a value belongs to exactly one variant among those
    listed (corresponding to tagged unions or \texttt{enums} in the code).
  \item \textbf{Vectors and Sequences} $\vec{X}$: Denotes a finite,
    ordered sequence of elements of type $X$ (corresponding to \texttt{Vec<X>}).
  \item \textbf{Optionality} $X^?$: Indicates a value that may be
    present (of type $X$) or absent (corresponding to \texttt{Option<X>}).
  \item \textbf{Function Signatures} $f : A \to B$: Denotes a
    mapping from domain $A$ to codomain $B$. Multi-argument functions
    use Cartesian products ($A \times B \to C$).
  \item \textbf{Boolean Domain} $\mathbb{B}$: The set of truth values
    $\{ \text{true}, \text{false} \}$.
\end{itemize}

\section{Glossary of Types and Domains}
The following definitions catalog the symbols that constitute the formal Intermediate Representation (IR) and the analysis pipeline.

\subsection{Base Types, Paths, and Identifiers}
\begin{itemize}
  \item $\mathcal{I}$ (\textbf{Identifier}): The set of valid strings
    used as names for variables, functions, types, and modules.
  \item $\Pi$ (\textbf{Path Domain}) and $\pi$ (\textbf{Qualified Name}):
    The domain of non-empty sequences of identifiers $\Pi = \mathcal{I}^+ =
    \{ \langle \iota_1, \dots, \iota_n \rangle \mid \iota_j \in \mathcal{I}, n \ge 1 \}$,
    where each $\pi \in \Pi$ uniquely identifies a component's absolute path.
  \item $\tau$ (\textbf{Type Reference}): An abstract reference to a type,
    represented by an eight-state sum type modeling its discovery lifecycle:
    \texttt{Prim}($\iota_p$) with $\iota_p \in \mathcal{P}$ (primitive keyword),
    \texttt{Unresolved}($\pi$) (raw path extracted from source),
    \texttt{ResolutionQuery}($q$) (algebraic navigation intent),
    \texttt{Resolved}($\pi$) (verified absolute path in the codebase),
    \texttt{External}($\pi$) (path resolved via external dependency),
    \texttt{Union}($\vec{\tau}$) (type disjunction or generic type arguments),
    \texttt{EvaluatedAccess}($\tau_{path}, \tau_{type}$) (preserves queried path and resolved target), or
    \texttt{Failed}($\pi$) (unresolved reference).
  \item $\delta$ (\textbf{Import Directive}): A dependency declaration
    $\langle \pi_{path}, \iota_{alias}^?, \mathbb{B}_{wildcard} \rangle$
    specifying an imported path, an optional local alias, and a wildcard flag.
\end{itemize}

\subsection{Structural Members and Signatures}
\begin{itemize}
  \item $v$ (\textbf{Variable}): A local variable declaration within a block,
    structurally equivalent to a field ($v \equiv f$).
  \item $f$ (\textbf{Field}): A member data field defined by the tuple
    $\langle \iota, \tau, \vec{\tau}_{ann} \rangle$.
  \item $p$ (\textbf{Parameter}): A formal parameter defined by
    $\langle \iota^?, \tau, \mathbb{B} \rangle$, comprising an optional
    name, a type reference, and a variadicity indicator.
  \item $\sigma$ (\textbf{Signature}): The functional signature tuple
    $\langle \vec{p}, \tau_{ret} \rangle$, decoupling inputs and output
    from the implementing body.
  \item $\mathcal{B}$ (\textbf{Block}): A lexical execution block
    $\langle \vec{f}_{local}, \vec{\tau}_{calls}, \vec{\tau}_{inst},
    \vec{\tau}_{accesses}, \vec{\tau}_{casts}, \vec{\mathcal{B}}_{sub} \rangle$.
  \item $\mathcal{F}$ (\textbf{Function}): A universal functional component
    $\langle \pi, \sigma, \mathcal{B}^?, \mathbb{B}_{ctor}, \vec{\tau}_{ann} \rangle$,
    binding a qualified name to a signature, an optional body block, a
    constructor flag, and annotations.
\end{itemize}

\subsection{Architectural Entities}
\begin{itemize}
  \item $\mathcal{S}$ (\textbf{Structured Type}): The abstraction for classes,
    structs, interfaces, traits, and enums, defined by
    $\langle \pi, k, \vec{f}, \vec{\mathcal{F}}, \vec{\tau}_{sup}, \vec{\mathcal{S}}_{nest}, \vec{\tau}_{ann}, \vec{\delta}_{imp} \rangle$
    where $k \in \{ \text{Class}, \text{Struct}, \text{Interface}, \text{Trait}, \text{Enum}, \text{EnumVariant} \}$.
  \item $\mathcal{A}$ (\textbf{Alias}): A type alias mapping a qualified
    name to an existing target type reference: $\langle \pi, \tau_{target} \rangle$.
  \item $\mathcal{X}$ (\textbf{Implementation Extension}): An implementation block
    aggregating methods and aliases for a target type $\tau_{target}$, flattened
    into $\mathcal{S}$ during Name Resolution.
  \item $\mathcal{C}$ (\textbf{Component}): The universal sum type of entities
    promoted to graph vertices: Modules ($\mathcal{M}$), Structured Types
    ($\mathcal{S}$), Type Aliases ($\mathcal{A}$), Functions ($\mathcal{F}$),
    Fields ($f$), Primitives, and External symbols.
  \item $\mathcal{M}$ (\textbf{Module}): The primary hierarchical namespace unit.
  \item $\mathcal{D}$ (\textbf{Intermediate Representation Domain}): The universe
    of valid module hierarchies $\vec{\mathcal{M}}$.
\end{itemize}

\subsection{Resolution Model and Scoping Contexts}
\begin{itemize}
  \item $\mathcal{Q}$ (\textbf{Query Algebra}): The domain of resolution expressions
    composed of primitives $\texttt{Find}(s)$, $\texttt{Extract}(q, s)$, and $\texttt{Call}(q)$, with $q \in \mathcal{Q}$.
  \item $\Sigma$ (\textbf{Scope}): A single environment of lexical visibility
    containing a symbol table, supertype inheritance links, and imports.
  \item $\mathcal{E}$ (\textbf{Scope Tree}): The workspace-wide hierarchical
    tree of $\Sigma$ scopes, enabling lexical scope climbing.
  \item $\Delta$ (\textbf{Symbol Stack}): The ephemeral stack of lexical frames
    $\Delta = [s_1, \dots, s_n]$ where each frame $s_i : \mathcal{I} \to \mathcal{Q}$
    maps local identifiers to queries during Query Building.
  \item $\mathcal{P}$ (\textbf{Primitive Registry}): The language-specific table
    of built-in primitive type keywords.
  \item $\mathcal{K}$ (\textbf{Language Configuration}): The configuration tuple
    governing language-specific semantic features (receiver keywords, impl blocks, forward declarations).
  \item $\Gamma$ (\textbf{Execution Context}): The immutable evaluation context
    $\Gamma = \langle \mathcal{E}, \mathcal{P}, \mathcal{K} \rangle$ used to resolve queries.
  \item $\mathcal{G}$ (\textbf{Dependency Graph}): The directed graph $\mathcal{G} = (V, E)$
    composed of component vertices $V$ and typed dependency edges $E$.
\end{itemize}

\subsection{Analysis Pipeline Phases}
\begin{itemize}
  \item $\phi$ (\textbf{Full Analysis Pipeline}): The end-to-end transformation
    $\phi : \mathcal{W} \to \mathcal{G}$, defined as $\phi = \gamma \circ \rho \circ \varepsilon$.
  \item $\varepsilon$ (\textbf{Workspace Extraction}): The language-dependent
    phase $\varepsilon : \mathcal{W} \to \vec{\mathcal{M}}$ translating source code into unresolved IR modules.
  \item $\rho$ (\textbf{Name Resolution}): The language-independent phase
    $\rho : \vec{\mathcal{M}} \to \vec{\mathcal{M}}$ that resolves symbolic type references into semantic pointers ($\rho = \rho_{\text{exec}} \circ \rho_{\text{build}}$).
    \begin{itemize}
      \item $\rho_{\text{build}}$ (\textbf{Query Building}): Traverses the IR with the Symbol Stack $\Delta$, transforming unresolved paths into algebraic queries $q \in \mathcal{Q}$.
      \item $\rho_{\text{exec}}$ (\textbf{Query Execution}): Evaluates queries $q$ against the Execution Context $\Gamma = \langle \mathcal{E}, \mathcal{P}, \mathcal{K} \rangle$.
    \end{itemize}
  \item $\gamma$ (\textbf{Graph Construction}): The final phase $\gamma : \vec{\mathcal{M}} \to \mathcal{G}$
    flattening the resolved IR into vertices $V$ and typed edges $E$.
\end{itemize}
